%% CV - Aditya Dasika, tailored for TCS: Gen AI Developer (Bengaluru)
%%
%% Compile with: lualatex main_tcs.tex
%% (pdflatex often fails on modern MiKTeX with fontawesome5 font-expansion errors.)

\documentclass[11pt,a4paper,sans]{moderncv}
\moderncvstyle{banking}
\moderncvcolor{blue}

\renewcommand*{\firstnamestyle}[1]{{\fontsize{34}{36}\bfseries\upshape\color{color1}#1}}
\renewcommand*{\lastnamestyle}[1]{{\fontsize{34}{36}\bfseries\upshape\color{color1}#1}}
\renewcommand*{\sectionstyle}[1]{{\sectionfont\color{color1}#1}}

\usepackage[utf8]{inputenc}
\usepackage{needspace}
\usepackage{hyperref}
\hypersetup{
    colorlinks=true,
    linkcolor=blue,
    filecolor=magenta,
    urlcolor=blue,
    pdftitle={Aditya Dasika - CV - TCS Gen AI Developer},
    pdfpagemode=FullScreen,
}
\usepackage[scale=0.80]{geometry}
\usepackage{import}

% personal data
\name{Aditya}{Dasika}
\address{Hyderabad, Telangana, India}{}{}
\phone[mobile]{+91-96184-33030}
\email{akhil.dasika47@gmail.com}
\extrainfo{\href{https://linkedin.com/in/aditya-dasika}{LinkedIn}, \href{https://github.com/aditya-dasika}{GitHub}}

\begin{document}

\makecvtitle

% ============================================================
%     PROFILE STATEMENT
% ============================================================

\vspace{6pt}
\small{Production AI/ML engineer specializing in agentic and generative AI systems. At Phenom People, architected chatbotCXAgent, a 12-node LangGraph orchestration platform handling intent recognition, dialogue state management, tool-routing, and Graph RAG retrieval, shipped to production users. Built the full agentic stack end to end: LangChain/LangGraph orchestration, FAISS-based RAG, OpenAI and Anthropic API integration, and XGrammar-enforced tool-calling reliability (91.9\% routing accuracy, zero malformed calls in production). Strong core Python fundamentals (OOP, decorators, typing, data structures) developed across 4+ years of production software engineering, including QLoRA fine-tuning on vLLM (10.98 RPS at 98.6\% success) and multi-task NLP with PyTorch. Built with Claude Code.}

% ============================================================
%     CORE COMPETENCIES
% ============================================================

\section{Core Competencies}
\vspace{1pt}
\begin{itemize}

\item \textbf{Generative \& Agentic AI}: LangGraph (production orchestration, LangGraph agents), LangChain, RAG and Graph RAG, FAISS vector search, OpenAI \& Anthropic APIs, tool/function calling, XGrammar constrained decoding

\item \textbf{Core Python}: OOP, decorators, typing/dataclasses, iterators/generators, exception handling, design patterns, data structures \& algorithms, applied across 4+ years of production systems

\item \textbf{ML \& NLP}: PyTorch, Transformers (BERT/ModernBERT), multi-task learning, NLU/NER/intent recognition, NumPy, Pandas, classical forecasting models

\item \textbf{LLM Fine-Tuning \& Serving}: QLoRA, SFT, Unsloth, HuggingFace PEFT, NF4 quantization, vLLM, Ollama

\item \textbf{Eval \& Observability}: pass@k, LLM-as-judge rubrics, LiteLLM, Langfuse, production MLOps

\item \textbf{Backend \& Cloud}: FastAPI, Node.js, Docker, AWS (EC2/S3/SageMaker), PostgreSQL, SQL, CI/CD; cloud ML pipeline patterns (training, serving, cost monitoring) built on AWS, transferable to GCP-equivalent services

\end{itemize}

% ============================================================
%     PROFESSIONAL EXPERIENCE
% ============================================================

\section{Professional Experience}
\vspace{3pt}
\begin{itemize}

\needspace{5\baselineskip}
\item{\cventry{Apr 2025--Present}{AI/LLM Engineer II}{Phenom People}{Hyderabad, Telangana}{}{\vspace{1pt}
\begin{itemize}
    \item Architected \textbf{chatbotCXAgent}, a production \textbf{agentic} platform on \textbf{LangGraph} (12 orchestrated nodes) handling intent recognition, dialogue state management, tool-routing, and \textbf{Graph RAG} retrieval across enterprise workflows, shipped to real users in production
    \item Enforced tool-calling reliability with XGrammar constrained decoding and schema-validated JSON outputs, raising tool-routing accuracy to \textbf{91.9\%} and eliminating malformed tool calls in production
    \item Fine-tuned and served Qwen models via QLoRA (Unsloth, NF4) on vLLM / AWS EC2 g5.xlarge (A10G), sustaining \textbf{10.98 RPS at 98.6\% success} under load and roughly halving inference cost vs.\ hosted GPT-4.1
    \item Built a custom eval framework (pass@k, LLM-judge rubric) gating every pipeline change across all 12 nodes before release; instrumented LiteLLM + Langfuse tracing for per-node cost and latency observability
    \item Fine-tuned a multi-task \textbf{PyTorch} ModernBERT model (intent classification + NER) improving NLU accuracy from \textbf{87\% to 97\%} via differential learning rates, label smoothing, and uncertainty-weighted multi-task loss
\end{itemize}}}

\vspace{3pt}

\needspace{5\baselineskip}
\item{\cventry{May 2022--Feb 2025}{Full Stack Engineer}{Blueleaves Farms}{Hyderabad, Telangana}{}{\vspace{1pt}
\begin{itemize}
    \item Built data-driven forecasting and analytics, including demand estimation, sales-trend, and inventory-planning models on transaction and stock-movement data, improving procurement and replenishment decisions
    \item Owned end-to-end delivery of production systems across POS, CRM, inventory, procurement, and e-commerce, serving \textbf{5+ locations and 10,000+ monthly transactions} while cutting manual operational errors by $\sim$90\%
    \item Designed scalable backend workflows on Firebase Cloud Functions, Node.js, and REST APIs, including invoicing, inventory synchronization, order lifecycle, and role-based access control, integrating multiple internal systems end to end
\end{itemize}}}

\end{itemize}

% ============================================================
%     PROJECTS
% ============================================================

\needspace{8\baselineskip}
\section{Projects}
\vspace{3pt}
\begin{itemize}

\needspace{4\baselineskip}
\item{\cventry{}{IndexNotes --- AI Knowledge Graph for Scientific Papers}{Neo4j, Ollama, PyMuPDF, Python, Docker}{}{}{\vspace{1pt}
\begin{itemize}
    \item Built an automated \textbf{RAG} pipeline ingesting scientific papers (PDF to local-LLM extraction to structured JSON to Neo4j), creating 40+ concepts and 17+ typed relationships across 3 research papers, designed for \textbf{Graph RAG} queries
\end{itemize}}}

\vspace{3pt}

\needspace{4\baselineskip}
\item{\cventry{}{Physics-Informed Neural Networks (PINNs)}{PyTorch, Autograd, L-BFGS, Latin Hypercube Sampling}{}{}{\vspace{1pt}
\begin{itemize}
    \item Implemented PINNs from scratch in \textbf{PyTorch} solving 3 PDEs: Heat equation ($\sim$3\% error), Burgers' equation (matching Raissi et al.), and Schr\"{o}dinger equation (complex-valued, final loss 0.000083)
\end{itemize}}}

\end{itemize}

% ============================================================
%     EDUCATION
% ============================================================

\section{Education}
\vspace{1pt}
\begin{itemize}

\item{\cventry{Aug 2018--May 2022}{Bachelor of Technology in Mechanical Engineering}{Andhra Loyola Institute of Engineering and Technology}{Vijayawada, Andhra Pradesh}{}{}}

\end{itemize}

% ============================================================
%     LANGUAGES
% ============================================================

\section{Languages}
\vspace{1pt}
\begin{itemize}
\item{English (fluent), Telugu (native)}
\end{itemize}

% ============================================================
%     REFERENCES
% ============================================================

\section{References}
\vspace{1pt}
\begin{itemize}
\item{Available upon request.}
\end{itemize}

\end{document}
